% GENERATED FILE. Edit canonical lessons, then run npm run generate:books.
% canonical-source-hash: fnv1a64:bc42769651a8137b
% canonical-lessons: PT-C359-gerente, PT-C359-departamento, PT-C359-tarefa, PT-C359-colaborador, PT-C359-delegar

\chapter{No escritório}
\label{ch:pt-no-escritorio}
\hlchaptermodality{359}

\begin{chapteropening}
\textbf{By the end of this chapter:} \emph{I can describe an office: the manager, the department, my tasks and colleagues, and what I delegate.}

Complete the last lesson of chapter 359: I can describe an office: the manager, the department, my tasks and colleagues, and what I delegate.
\end{chapteropening}

\section[o gerente]{o gerente — the manager}

\label{lesson:PT-C359-gerente}



\begin{quote}
Before the new one: say the Portuguese for the strike, then the Portuguese for to lay claim to.
\end{quote}

\subsection*{You'll want to know: o gerente}
\textbf{o gerente} — \textquotedblleft{}the manager\textquotedblright{}.

\begin{quote}
O gerente da loja resolveu o problema.
\end{quote}

\begin{quote}
\emph{The shop manager sorted out the problem.}
\end{quote}

Same form for a woman: \emph{a gerente}.

\begin{cousinweb}
\emph{gerente} is a learned borrowing from Latin \emph{gerēns} (EtymDB), the present participle of \emph{gerere} \textquotedblleft{}to carry, to carry on, to manage\textquotedblright{} (Lewis \& Short).
\end{cousinweb}

\begin{grammarlens}[title={What you've built}]
One more word for telling, explaining and describing at B1.
\end{grammarlens}

\subsection*{Your turn}
\begin{itemize}
\raggedright
  \item \emph{Say it:} \emph{o gerente}
  \item \emph{Say it:} \emph{o gerente}, once more
  \item \emph{Say it:} a sentence of your own with \emph{o gerente}
  \item \emph{Recall:} say the Portuguese for the strike, then the Portuguese for to lay claim to, then say \emph{o gerente} again
\end{itemize}

\begin{tcolorbox}[breakable,colback=teal!4,colframe=teal!35!black,title={Before you move on}]
What does o gerente mean? (\textquotedblleft{}The manager\textquotedblright{}.) Say it once more.
\end{tcolorbox}
\section[o departamento]{o departamento — the department}

\label{lesson:PT-C359-departamento}



\begin{quote}
Before the new one: say the Portuguese for to lay claim to, then the Portuguese for the manager.
\end{quote}

\subsection*{You'll want to know: o departamento}
\textbf{o departamento} — \textquotedblleft{}the department\textquotedblright{}.

\begin{quote}
Ela trabalha no departamento de vendas.
\end{quote}

\begin{quote}
\emph{She works in the sales department.}
\end{quote}

\emph{O departamento de recursos humanos} — the human resources department.

\begin{cousinweb}
\emph{departamento} was borrowed from French \emph{département} (EtymDB), built on Old French \emph{departir} \textquotedblleft{}to divide\textquotedblright{}, which goes back to Latin \emph{partīre} \textquotedblleft{}to divide\textquotedblright{}, on \emph{pars} \textquotedblleft{}part\textquotedblright{} (Lewis \& Short).
\end{cousinweb}

\begin{grammarlens}[title={What you've built}]
One more word for telling, explaining and describing at B1.
\end{grammarlens}

\subsection*{Your turn}
\begin{itemize}
\raggedright
  \item \emph{Say it:} \emph{o departamento}
  \item \emph{Say it:} \emph{o departamento}, once more
  \item \emph{Say it:} a sentence of your own with \emph{o departamento}
  \item \emph{Recall:} say the Portuguese for to lay claim to, then the Portuguese for the manager, then say \emph{o departamento} again
\end{itemize}

\begin{tcolorbox}[breakable,colback=teal!4,colframe=teal!35!black,title={Before you move on}]
What does o departamento mean? (\textquotedblleft{}The department\textquotedblright{}.) Say it once more.
\end{tcolorbox}
\section[a tarefa]{a tarefa — the task, the job (to be done)}

\label{lesson:PT-C359-tarefa}



\begin{quote}
Before the new one: say the Portuguese for the manager, then the Portuguese for the department.
\end{quote}

\subsection*{You'll want to know: a tarefa}
\textbf{a tarefa} — \textquotedblleft{}the task, the job (to be done)\textquotedblright{}.

\begin{quote}
Hoje tenho três tarefas importantes.
\end{quote}

\begin{quote}
\emph{Today I have three important tasks.}
\end{quote}

\emph{Tarefas domésticas} — housework.

\begin{cousinweb}
\emph{tarefa} was borrowed from Arabic \emph{ṭarīḥa} \textquotedblleft{}endeavour, a task set\textquotedblright{} (EtymDB); Spanish \emph{tarea} has the same source.
\end{cousinweb}

\begin{grammarlens}[title={What you've built}]
One more word for telling, explaining and describing at B1.
\end{grammarlens}

\subsection*{Your turn}
\begin{itemize}
\raggedright
  \item \emph{Say it:} \emph{a tarefa}
  \item \emph{Say it:} \emph{a tarefa}, once more
  \item \emph{Say it:} a sentence of your own with \emph{a tarefa}
  \item \emph{Recall:} say the Portuguese for the manager, then the Portuguese for the department, then say \emph{a tarefa} again
\end{itemize}

\begin{tcolorbox}[breakable,colback=teal!4,colframe=teal!35!black,title={Before you move on}]
What does a tarefa mean? (\textquotedblleft{}The task, the job (to be done)\textquotedblright{}.) Say it once more.
\end{tcolorbox}
\section[o colaborador]{o colaborador — the staff member, the colleague; the collaborator}

\label{lesson:PT-C359-colaborador}



\begin{quote}
Before the new one: say the Portuguese for the department, then the Portuguese for the task.
\end{quote}

\subsection*{You'll want to know: o colaborador}
\textbf{o colaborador} — \textquotedblleft{}the staff member, the colleague; the collaborator\textquotedblright{}.

\begin{quote}
A empresa tem cem colaboradores.
\end{quote}

\begin{quote}
\emph{The company has a hundred staff.}
\end{quote}

Companies in Portugal often call their staff \emph{colaboradores}. Feminine \emph{a colaboradora}.

\begin{cousinweb}
\emph{colaborador} is built on \emph{colaborar} \textquotedblleft{}to work together\textquotedblright{}, a learned borrowing from Late Latin \emph{collabōrāre}, from \emph{con-} and \emph{labōrāre} \textquotedblleft{}to work\textquotedblright{} (Lewis \& Short).
\end{cousinweb}

\begin{grammarlens}[title={What you've built}]
One more word for telling, explaining and describing at B1.
\end{grammarlens}

\subsection*{Your turn}
\begin{itemize}
\raggedright
  \item \emph{Say it:} \emph{o colaborador}
  \item \emph{Say it:} \emph{o colaborador}, once more
  \item \emph{Say it:} a sentence of your own with \emph{o colaborador}
  \item \emph{Recall:} say the Portuguese for the department, then the Portuguese for the task, then say \emph{o colaborador} again
\end{itemize}

\begin{tcolorbox}[breakable,colback=teal!4,colframe=teal!35!black,title={Before you move on}]
What does o colaborador mean? (\textquotedblleft{}The staff member, the colleague; the collaborator\textquotedblright{}.) Say it once more.
\end{tcolorbox}
\section[delegar]{delegar — to delegate}

\label{lesson:PT-C359-delegar}



\begin{quote}
Before the new one: say the Portuguese for the task, then the Portuguese for the staff member.
\end{quote}

\subsection*{You'll want to know: delegar}
\textbf{delegar} — \textquotedblleft{}to delegate\textquotedblright{}.

\begin{quote}
A chefe delegou o trabalho em dois colegas.
\end{quote}

\begin{quote}
\emph{The boss delegated the work to two colleagues.}
\end{quote}

\emph{Delegar em alguém} — to delegate to someone.

\begin{cousinweb}
\emph{delegar} is a learned borrowing from Latin \emph{dēlēgāre} \textquotedblleft{}to send, to assign, to delegate\textquotedblright{} (EtymDB, Lewis \& Short), built from \emph{dē} and \emph{lēgāre} \textquotedblleft{}to send, to depute\textquotedblright{}.
\end{cousinweb}

\begin{grammarlens}[title={What you've built}]
One more word for telling, explaining and describing at B1.
\end{grammarlens}

\subsection*{Your turn}
\begin{itemize}
\raggedright
  \item \emph{Say it:} \emph{delegar}
  \item \emph{Say it:} \emph{delegar}, once more
  \item \emph{Say it:} a sentence of your own with \emph{delegar}
  \item \emph{Recall:} say the Portuguese for the task, then the Portuguese for the staff member, then say \emph{delegar} again
\end{itemize}

\begin{tcolorbox}[breakable,colback=teal!4,colframe=teal!35!black,title={Before you move on}]
What does delegar mean? (\textquotedblleft{}To delegate\textquotedblright{}.) Say it once more.
\end{tcolorbox}
